\documentclass[12pt]{article}
\usepackage[utf8]{inputenc}
\usepackage[T2A]{fontenc}
\usepackage[english,russian]{babel}
\usepackage{amsmath}
\usepackage{amssymb}
\usepackage[a4paper,margin=2cm]{geometry}
\usepackage{setspace}
\usepackage{hyperref}
\usepackage{mathtools}
\onehalfspacing
\sloppy

\title{Общее решение. Общий интеграл. Общее решение в параметрической форме.}

\begin{document}

\maketitle

\section{Общее решение}
Рассмотрим семейство решений дифференциального уравнения $y' = f(x, y)$, зависящее от одного параметра $C$:
\[y = \varphi(x, C).\]
Такое семейство решений называется \textbf{общим решением} дифференциального уравнения.

Геометрически общее решение представляет собой совокупность всех интегральных кривых данного уравнения
зависящих от параметра $C$.

При помощи общего решение можно найти решение задачи Коши, подставив в
общее решение значение $x_0$ и $y_0$ из условия задачи Коши и найдя соответствующее
значение параметра $C$. При этом единственность и разрешимость задачи Коши при таком
подходе не гарантируется.

\section{Общее решение в области $D$ изменения переменных $x$ и $y$}
Пусть $D$ - облатсь на плоскости $Oxy$, через каждую точку которой
проходит ровно одна интегральная кривая дифференциального уравнения.
В каждой точке области $D$ имеет место существование и единственность решения задачи Коши.



\end{document}